\section{SCENARIO \#14: "The Woodstock Races'': Oct. 6-11, 1864
}


\begin{scenariodata}
\textbf{Game Length:} 6 Turns\\
\end{scenariodata}

\begin{scenariodata}
(Weather: July - October Column)\\
\end{scenariodata}

\begin{scenariodata}
\textbf{First Phase:} Union\\
\end{scenariodata}

\begin{scenariodata}
\textbf{Union:}\\
\end{scenariodata}

\begin{scenariodata}
Resource Factor 217\\
\end{scenariodata}

\begin{scenariodata}
\textbf{Initial Placement:}\\
\end{scenariodata}

\begin{scenariodata}
Hex A-16 (Romney) : (41, 1A], (1F, 1S)\\
\end{scenariodata}

\begin{scenariodata}
Hex 0-44 :\\
\end{scenariodata}

\begin{scenariodata}
\textbf{and adjacent :} 13C, 4H\\
\end{scenariodata}

\begin{scenariodata}
Hex 0-42 (Harrisonburg) : 651, 9A,\\
\end{scenariodata}

\begin{scenariodata}
\textbf{and adjacent :} (Wright, Emory, 1S, 2W)\\
Hex M-32 > 121, 3A,\\
\end{scenariodata}

\begin{scenariodata}
\textbf{and adjacent :} (Crook, 1S, 1W)\\
\end{scenariodata}

\begin{scenariodata}
Hex M-29 (Woodstock) : 21, 1C, (Sheridan)\\
\end{scenariodata}

\begin{scenariodata}
Hex N-20 (Winchester) : 41, (1F, 5S, 3W)\\
\end{scenariodata}

\begin{scenariodata}
Hex S-28 (Front Royal\} : 2f, (1W)\\
\end{scenariodata}

\begin{scenariodata}
Hex M-26 (Strasburg) : 31, (1F, 1S, 1W)\\
\end{scenariodata}

\begin{scenariodata}
Hex R-46 (Port Republic) : 6C, 2H, (1S, 1W)\\
\end{scenariodata}

\begin{scenariodata}
B\&O + [201, 5A, 2C], (4F, 45)\\
\end{scenariodata}

\begin{scenariodata}
\textbf{Withdrawal Hexes:}\\
\end{scenariodata}

\begin{scenariodata}
N-7, 8-10, Z-13, CC-15 through CC-33\\
\end{scenariodata}

\begin{scenariodata}
(Supply/Wagon Entry Hexes:\\
\end{scenariodata}

\begin{scenariodata}
N-7, S-10, Z-13, A-11, A-13, A-16, CC-15 through CC-33)\\
\end{scenariodata}

\begin{scenariodata}
No additional Fortification Centers.\\
\end{scenariodata}

\begin{scenariodata}
\textbf{Confederates:}\\
\end{scenariodata}

\begin{scenariodata}
Resource Factor 714\\
\end{scenariodata}

\begin{scenariodata}
\textbf{Initial Placement:}\\
\end{scenariodata}

\begin{scenariodata}
Hex P-48 : 461, 11A, 13C, 4H\\
\end{scenariodata}

\begin{scenariodata}
\textbf{and adjacent :} (Early, Gordon, Ramseur,\\
\end{scenariodata}

Kershaw, 4F, 3S, 6W)

\begin{scenariodata}
Hex 0-51 (Staunton) > 11, (1F)\\
\end{scenariodata}

\begin{scenariodata}
Hex BB-53 (Charlottesville) : 21, 1A, (1S, 1F\}\\
\end{scenariodata}

\begin{scenariodata}
(Partisans: Mosby, Gilmor)\\
\end{scenariodata}

\begin{scenariodata}
\textbf{Withdrawal Hexes:}\\
\end{scenariodata}

\begin{scenariodata}
CC-40 through CC-58\\
\end{scenariodata}

\begin{scenariodata}
\textbf{Supply/Wagon Entry Hexes:}\\
\end{scenariodata}

\begin{scenariodata}
CC-40 through CC-58, 0-51, M-50, E-46\\
\end{scenariodata}

\begin{scenariodata}
No additional Fortification centers.\\
\end{scenariodata}

\begin{scenariodata}
SPECIAL \#1: Hex O-51 (Staunton) is worth no Victory Points to the\\
Union in this scenario.\\
\end{scenariodata}

\begin{scenariodata}
(SPECIAL \#2: Place a Destruction marker on the bridge at O-44/0-45.)\\
\end{scenariodata}

\begin{scenariodata}
SPECIAL \#3: Devastated areas are shown by placing Destruction\\
\end{scenariodata}

counters in the following hexes: R-23, P-24, 0-21, M-22, L-19, M-49,

\begin{scenariodata}
N-47, P-50, Q-48, T-50, S-52.)\\
\end{scenariodata}

\begin{scenariodata}
\textbf{Victory Conditions:} All other results are a draw.\\
\end{scenariodata}

\begin{scenariodata}
\textbf{ADVANCED GAME:} Union wins if they have 20 or more lead in total\\
\end{scenariodata}

\subsection*{Victory Points.}

\begin{scenariodata}
Confederates win if they have 15 or more lead in\\
\end{scenariodata}

total Victory Points.

\begin{scenariodata}
\textbf{BASIC GAME:} Same as the Advanced Game.\\
\end{scenariodata}

\begin{scenariodata}
\textbf{HISTORICAL RESULTS:} Union victory.\\
\end{scenariodata}

\begin{scenariodata}
\textbf{PLAYTESTER'S COMMENTS:} The Union Player can win by withdraw-\\
\end{scenariodata}

ing to the north, devastating as he goes, but such a strategy will force the army into a large number of small Stacks, that the Confederates can fall on and defeat in detail. A Union cavalry screen is useful in preventing this. An alternate Union strategy would be to push for Charlottesville and try to force a battle to the west of the Blue Ridge. The Confederates should watch for an opening, and try to crush some isolated Union force.

| eee |

Sheridan's Federals systematically devastated the Shenandoah while retiring. Early main Rebel army followed at some distance, while the two cavalry forces clashed constantly. A mass Union cavalry charge at Woodstock ended the contest, and the victorious Yankee cavalry chased the beaten Southern horsemen for twenty-six miles. The Union army completed its maneuvers without further molestation, and took up a strong position on Cedar Creek.

Sheridan regarded the campaign as virtually over. He had taken a hastily-thrown together assortment of units, molded them into an effective army, and led them to a series of triumphs that made them the heroes of the North. The burned-out Valley of Virginia had lost much of its former value to the Confederacy; and, Jubal Early, once hailed throughout the South as another 'Stonewall,' was largely discredited, keeping his command only through the continued support of Robert E. Lee. Early's army appeared to be finished, and even the dreaded partisans seemed to be on the ropes (Mosby and Gilmor had both been wounded during the campaign, and McNeil had been killed). Sheridan himself went north to confer on the further strategy proposed for his army. While he was still absent, the Confederates decided to make one last, desperate bid for victory. Early moved his army up to the Union positions, sent Gordon and the Il Corps on an all-night flank march, and launched an attack at dawn...

